\chapter{Tools and engineering process}
This chapter provides an overview of the tools and processes used throughout the development of Miwo. While the background chapter lays the theoretical groundwork, this chapter focuses on the practical approaches for designing, building, and evaluating the robot. This includes the physical work, such as 3D printing, electronics, and mechanical assembly, as well as the software frameworks and tools used to develop, test, and validate the system.
Table \ref{tab:tools} shows an overview of the tools and software used:

\begin{table}[h]
    \centering
    \begin{tabular}{|l|l|l|}
        \hline
        \textbf{Tool} & \textbf{Name}& \textbf{Version} \\
        \hline
        & & \\[1em]
        \hline
        & & \\[1em]
        \hline
        & & \\[1em]
        \hline
        & & \\[1em]
        \hline
        & & \\[1em]
        \hline
        & & \\[1em]
        \hline
        & & \\[1em]
        \hline
        & & \\[1em]
        \hline
        & & \\[1em]
        \hline
        & & \\[1em]
        \hline
    \end{tabular}
    \caption{Table showing the tools and software used in the work on the thesis.}
    \label{tab:tools}
\end{table}

\newpage
\section*{ROS2}

\section*{3D design (CAD)}
Autodesk Fusion (formerly known as Fusion 360) is a cloud-based platform that combines computer-aided design (CAD), CAM, CAE, and PCB software development into one single tool. It offers a wide range of tools for 3D modeling, design, simulation, and engineering, making it a versatile tool for professionals, researchers, and students. Autodesk Fusion also supports parametric modeling, meaning that design geometry is created using parameters, constraints, and features. This allows designers to change a model's shape instantly by updating it rather than redrawing, ensuring that changes to one part automatically propagate through the entire assembly. All design work for Miwo was carried out in Autodesk Fusion, including structural components, support structures, and body parts that were subsequently 3D printed.

\section*{3D Printing}
3D printing, or additive manufacturing, is a process that allows the creation of three-dimensional objects from digital files by depositing materials layer by layer. 3D printing has grown rapidly in popularity among researchers, professionals, and hobbyists for its potential to enable fast, affordable prototyping across a variety of industries.
During an engineering process, 3D printing allows users to create complex geometries and parts from a wide variety of materials while remaining cost-effective and rapid. With growing popularity and interest, there is a never-ending array of printers, ranging from cheap to expensive, for more detailed and precise work, making them widely accessible.

\section{Electronics and mechanics}

\subsection*{Onboard computer}
The NVIDIA Jetson Orin Nano developer kit is a compact, powerful single-board computer platform from NVIDIA, designed for applications such as generative AI, vision AI, and robotics. It features a powerful ARM-based CPU with an integrated NVIDIA GPU, enabling it to run complex computations without relying on an external power supply. The development board has a small form factor and low power consumption, making it well-suited for portable and embedded robotic systems where weight and size may be constrained. The affordable price point also provides developers, students, and researchers with a powerful and accessible platform. The Jetson Orin nano serves as a suitable onboard computer for mobile robot applications by providing sufficient computing power to run ROS2 nodes and handle communication between the platform's components via USB serial connections.
\subsection*{Servo motors}
Servo motors are rotary actuators that allow for precise control of angular position, velocity, and acceleration. There is a wide variety of servo motors available, ranging from low-cost to high-end professional actuators. Servo motors are used for many robotic applications, such as precise joint control with position feedback or simpler PWM-based alternatives.
\subsection{Go Robot actuators}
\subsection{Power supply}
\subsection{IMU}

\section{Simulation tools and environments}
\subsection{Rviz}
Rviz is a ROS2-compatible 3D visualization tool that allows developers to inspect and debug robotic systems in real time. It provides a user-friendly graphical interface for displaying relevant data, including coordinate frames, joint states, robot models, and sensor data, making it a powerful tool for visualizing the development and testing of robot systems. Rviz works by subscribing to data published to ROS" topics and then rendering them visually, allowing the developer to observe the state of the robot without needing to interact directly with the physical hardware, providing a more intuitive experience.
\subsection{Gazebo}
Gazebo is a powerful open-source 3D robotics simulator compatible with ROS2. Gazebo differs from other pure visualization tools, like Rviz, by providing a powerful physics engine that simulates robots and their interactions with the environment. By using Gazebo, we simulate physical behavior from the ground up, allowing us to model complex physics such as gravity, contact forces, and inertia. This allows the developer to test control algorithms and motion sequences using simulated physics before deploying them on physical hardware.
\subsection{Isaac Sim}
"Isaac Sim is an open source reference framework developed by NVIDIA, built on Omniverse libraries for robotics simulation, testing, and synthetic data generation in physically based virtual environments." Isaac Sim offers physically accurate simulated environments with GPU-accelerated physics and rendering, making it well-suited for applications like reinforcement learning, where large numbers of simulation steps are computed in a short time. Unlike Gazebo, which runs on the CPU, Isaac Sim leverages the
processing capabilities of modern GPUs to run thousands of simulations 
instances almost simultaneously, dramatically accelerating the training process.



